%  proposal.tex  --  MATH 6660, Final Project proposal (due Mon Nov 9, 2026).
%  Build: pdflatex proposal.tex  (twice)
\documentclass[12pt]{article}
\usepackage[letterpaper,margin=1in]{geometry}
\usepackage[T1]{fontenc}
\usepackage{amsmath}
\usepackage{newtxtext,newtxmath}   % Times New Roman text and math
\usepackage{bm}
\usepackage{microtype}
\usepackage{enumitem}
\usepackage{xcolor}
\usepackage[colorlinks=true,linkcolor={blue!45!black},citecolor={blue!45!black},
            urlcolor={blue!45!black},pdfborder={0 0 0}]{hyperref}

\setlength{\parindent}{0pt}
\setlength{\parskip}{6pt}
\newcommand{\jump}[1]{[\![#1]\!]}

\hypersetup{pdftitle={MATH 6660 Final Project Proposal},
            pdfauthor={Md Mahmudul Hasan Mollah}}

\begin{document}
\thispagestyle{empty}

\begin{center}
  {\large\bfseries MATH 6660 --- Final Project Proposal}\\[2pt]
  {\small Fall 2026}\\[8pt]
  Md Mahmudul Hasan Mollah\\
\end{center}

\vspace{4pt}
\hrule
\vspace{10pt}

\textbf{Title.}
Cohesive Zone Elements: Camanho's Bilinear Law,
its Finite Element Implementation, and its Engineering Applications

\textbf{Article.}
P.~P.~Camanho, C.~G.~D\'avila, and M.~F.~de~Moura,
``Numerical simulation of mixed-mode progressive delamination in composite
materials,'' \emph{Journal of Composite Materials} \textbf{37}(16) (2003),
1415--1438. \href{https://doi.org/10.1177/0021998303034505}{doi:10.1177/0021998303034505}

\textbf{Proposed project.}
Camanho, D\'avila and de~Moura model interface fracture with zero-thickness
decohesion elements governed by a bilinear mixed-mode traction--separation
law: a quadratic stress criterion for damage onset, the Benzeggagh--Kenane
criterion for propagation, and a single irreversible damage
variable~\cite{camanho2003,camanho2002}.
This project will derive the finite element formulation, in which the
displacement may jump across the fracture surface $\Gamma_c$ and the weak form
gains an interface term, and show how the damage evolves under different
loading conditions (pure opening, pure shear and compression--shear) using
single-element simulations in MOOSE~\cite{moose}, an open-source finite element
framework that implements this law. Finally, the project will apply the model to example problems, which may include composite materials, nuclear reactor components and natural fractures in subsurface rock.\\

\vspace{4pt}
\begingroup
\small
\renewcommand{\section}[2]{\textbf{References}\par\vspace{-2pt}}
\begin{thebibliography}{9}
  \setlength{\itemsep}{1pt}
  \bibitem{camanho2003}
    P.~P.~Camanho, C.~G.~D\'avila, M.~F.~de~Moura,
    Numerical simulation of mixed-mode progressive delamination in composite
    materials, \emph{J.\ Compos.\ Mater.} 37 (2003) 1415--1438.
  \bibitem{camanho2002}
    P.~P.~Camanho, C.~G.~D\'avila,
    Mixed-mode decohesion finite elements for the simulation of delamination
    in composite materials, NASA/TM-2002-211737, NASA Langley Research
    Center, 2002.
  \bibitem{moose}
    G.~Giudicelli et al., 3.0 -- MOOSE: Enabling massively parallel
    multiphysics simulations, \emph{SoftwareX} 26 (2024) 101690.
\end{thebibliography}
\endgroup

\end{document}
